% TOP_PROBLEM: {"id": 3157, "title": "Subgraph of large average degree and large girth.", "category_id": 3, "category": {"id": 3, "name": "graph_theory", "display_name": "Graph Theory", "description": "Problems involving graphs, networks, and their properties.", "slug": "graph-theory", "order_index": 3, "created_at": "2026-07-31T15:26:25.671Z"}, "status": "open", "problem_number": "OPG-46708", "set_id": 12}
% FIRST_POSTED: 2026-10-01
% ENTRY_KIND: statement_only
% UnsolvedMath ID 3157; source catalogue code OPG-46708
% !TeX program = lualatex
% Definition and sources only; this file is not an LLM solution attempt.
\documentclass[11pt,a4paper]{article}
\usepackage[margin=25mm]{geometry}
\usepackage{fontspec}
\setmainfont{lmroman10-regular.otf}[BoldFont=lmroman10-bold.otf,
 ItalicFont=lmroman10-italic.otf,BoldItalicFont=lmroman10-bolditalic.otf]
\usepackage{amsmath,amssymb,mathtools}
\usepackage{unicode-math}
\setmathfont{latinmodern-math.otf}
\AtBeginDocument{\let\setminus\smallsetminus}
\usepackage{xurl,hyperref}
\hypersetup{colorlinks=true,urlcolor=blue}
\setcounter{secnumdepth}{0}
\setlength{\parindent}{0pt}
\setlength{\parskip}{5pt}
\setlength{\emergencystretch}{3em}
\begin{document}
\section*{Subgraph of large average degree and large girth.}\label{3157}
{\small UnsolvedMath ID 3157 · OPG-46708 · Graph Theory · OpenGarden}
\subsection{Definitions and mathematical statement}
For a nonempty finite simple graph $G=(V,E)$, its average degree
is $\overline d(G)=2|E|/|V|$. Its girth is the minimum length of
a simple cycle; set the girth to $+\infty$ when the graph is a
forest. A subgraph may be obtained by deleting vertices and
edges, so it need not be induced or spanning.

The conjecture has the following quantifier order:
\[
 \forall g,k\in\mathbb Z_{\geq1}\quad
 \exists d=d(g,k)\in\mathbb Z_{\geq1}\quad
 \forall G\quad
 \overline d(G)\geq d\Longrightarrow
 \exists H\subseteq G:
 \overline d(H)\geq k,\quad\operatorname{girth}(H)>g.
\]
Here $G$ and the selected $H$ are nonempty finite simple graphs.
The threshold is independent of the order and maximum degree
of $G$. The conclusion means that $H$ contains no cycle with at
most $g$ edges while retaining the prescribed average degree.

No regularity assumption is imposed on $G$. In particular, its
large average degree may coexist with widely varying vertex
degrees. One may discard edges of short cycles and choose a
smaller vertex set, but both resulting conditions must hold in
the same subgraph. The statement demands an existence threshold
for each pair $(g,k)$; it does not prescribe a numerical formula
or a polynomial bound for $d(g,k)$. A counterexample would fix
some pair $(g,k)$ and produce graphs of unbounded average degree
with no suitable subgraph.
\subsection{Short English statement}
Does sufficiently large average degree force a subgraph with any prescribed average degree and no short cycles?
\subsection{Sources}
\begin{itemize}
\item[S1] ULAM.AI and the UnsolvedMath Contributors, supplied problems.json, record 3157 (OPG-46708), OpenGarden collection. Catalogue identity and source transcription; CC BY 4.0.
\url{https://huggingface.co/datasets/ulamai/UnsolvedMath}.
\item[S2] Open Problem Garden, Subgraph of large average degree and large girth.. Original problem and discussion; the mathematical formulation below is expanded from the supplied source record.
\url{http://www.openproblemgarden.org/op/subgraph_of_large_average_degree_and_large_average_degree}.
\end{itemize}
\end{document}
